
\paragraph{\textbf{Corollary 6.4.}}
\SGARefTarget{sga3:xiv:6.4}
{\itshape Let \(G\) be a smooth algebraic group over an infinite
field \(k\). Then the set of \(k\)-rational points of
\(\mathcal C\), in the notation of \SGARefLink{sga3:xiv:6.1}{6.1}, is Zariski dense. The
union of the Cartan subgroups of \(G\) is dense in \(G\).}

The first assertion holds for every rational variety over an
infinite field and is, for us, the most important arithmetic
consequence of the rationality result. The second follows from
the first and the density theorem \SGARefLink{sga3:xiii:theorem:2:1}{XIII~2.1}.

\paragraph{\textbf{Corollary 6.5.}}
\SGARefTarget{sga3:xiv:6.5}
{\itshape Let \(G\) be a smooth connected algebraic group over
\(k\). Then the variety \(Z\) of regular semisimple points of
\(G\) \textup{(\SGARefLink{sga3:xiii:corollary:3:5}{XIII~3.5})} is unirational. In particular, if
\(k\) is infinite, its \(k\)-rational points are dense in \(Z\).}

Indeed, \(Z\) is an open subscheme of a torus over
\(\mathcal C\). Its function field \(L\) is therefore the
function field of a torus over the function field \(K\) of
\(\mathcal C\), namely the generic maximal torus of \(G\). Thus
\(L\) is unirational over \(K\) by \SGARefLink{sga3:xiii:lemma:3:4}{XIII~3.4}. Since \SGARefLink{sga3:xiv:6.1}{6.1} makes
\(K\) purely transcendental over \(k\), \(L\) is unirational
over \(k\).

\paragraph{\textbf{Corollary 6.6.}}
\SGARefTarget{sga3:xiv:6.6}
{\itshape Let \(G\) be a smooth connected algebraic group over
\(k\), and let \(H\) be the subgroup generated by the subvariety
\(Z\) of regular semisimple points. Equivalently
\textup{(\SGARefLink{sga3:xii:8.2}{XII~8.2})}, \(H\) is the smallest invariant algebraic
subgroup of \(G\) for which \(G/H\) has reductive rank zero; over
\(\bar k\), \(G/H\) is an extension of an abelian variety by a
smooth connected unipotent group. Then \(H\) is unirational.

In particular, suppose \(G=H\); equivalently
\textup{(\SGARefLink{sga3:xii:8.4}{XII~8.4})}, suppose \(G\) is affine and there is no
nontrivial homomorphism \(G_{\bar k}\to\mathbb G_a\). Then \(G\)
is unirational. If, in addition, \(k\) is infinite, \(G(k)\) is
dense in \(G\).}

This follows immediately from \SGARefLink{sga3:xiv:6.5}{6.5}. Indeed, if \(Z_i\) are smooth
connected unirational \(k\)-preschemes and
\(u_i\colon Z_i\to G\) are morphisms, then the algebraic subgroup
of \(G\) generated by the \(u_i\) is unirational
\textup{(VI\textsubscript{B}~\SGARefLink{prop:sga3-VIB-7.1}{7.1})}.

An interesting special case of either \SGARefLink{sga3:xiv:6.5}{6.5} or \SGARefLink{sga3:xiv:6.6}{6.6} is the
following.

\paragraph{\textbf{Corollary 6.7.}}
\SGARefTarget{sga3:xiv:6.7}
{\itshape Let \(G\) be a smooth connected affine algebraic group
of unipotent rank zero. Then \(G\) is unirational.}

The conclusion of \SGARefLink{sga3:xiv:6.4}{6.4} may be sharpened as follows.

\paragraph{\textbf{Corollary 6.8.}}
\SGARefTarget{sga3:xiv:6.8}
{\itshape Let \(G\) be a smooth connected algebraic group over an
infinite field \(k\), and let \(C\) be a Cartan subgroup of \(G\).
Then the union of the conjugates of \(C\) by regular semisimple
elements of \(G(k)\) is dense in \(G\).}

This follows at once from \SGARefLink{sga3:xiv:6.4}{6.4} and \SGARefLink{sga3:xiii:corollary:3:6}{XIII~3.6}, which says that the
morphism
\[
  \varphi\colon Z\times C\longrightarrow G,\qquad
  \varphi(t,c)=\operatorname{Ad}(t)c,
\]
is dominant. The same result also implies the next corollary,
without any assumption that \(k\) is infinite.

\paragraph{\textbf{Corollary 6.9.}}
\SGARefTarget{sga3:xiv:6.9}
{\itshape Let \(G\) be a smooth connected algebraic group over
\(k\), and let \(H\) be a smooth connected algebraic subgroup
having the same reductive rank and the same nilpotent rank as
\(G\). Equivalently, over \(\bar k\), the group \(H_{\bar k}\)
contains a Cartan subgroup of \(G_{\bar k}\). If \(H\) is
unirational, so is \(G\). If \(H(k)\) is dense in \(H\), then
\(G(k)\) is dense in \(G\).}

The morphism
\[
  \varphi\colon Z\times H\longrightarrow G,\qquad
  \varphi(t,h)=\operatorname{Ad}(t)h,
\]
is dominant. By \SGARefLink{sga3:xiv:6.5}{6.5}, \(Z\) is unirational, and by hypothesis so
is \(H\), hence so are \(Z\times H\) and \(G\). The density
assertion is proved in the same way.

We recover the following well-known result, due to Chevalley in
characteristic zero and to Rosenlicht in characteristic
\(p>0\).

\paragraph{\textbf{Corollary 6.10.}}
\SGARefTarget{sga3:xiv:6.10}
{\itshape Let \(G\) be a smooth connected affine algebraic group
over a perfect field \(k\). Then \(G\) is unirational. If \(k\)
is also infinite, \(G(k)\) is dense in \(G\).}

By \SGARefLink{sga3:xiv:1.1}{1.1}, \(G\) has a Cartan subgroup \(C\). By \SGARefLink{sga3:xiv:6.9}{6.9}, it is enough
to prove that \(C\) is unirational. Since \(k\) is perfect,
Galois descent from the algebraically closed case
\textup{\emph{Bible}~6, \SGARefLink{sga3:xiv:7:lemma2}{Theorem~2}} gives
\[
  C=T\times C_u,
\]
where \(T\) is the maximal torus of \(C\) and \(C_u\) is smooth,
connected, and unipotent. The torus \(T\) is unirational by
XIII~3.4. Moreover, \(C_u\) is \(k\)-unipotent because \(k\) is
perfect, so \SGARefLink{sga3:xiv:6.3}{6.3} applies to \(C_u\).
